\clearpage
\section{Nạp lại bản đồ và định vị bằng AMCL}
\subsection{Nguồn bản đồ và cách map server diễn giải ảnh}
\code{maps/map.pgm} là bản sao ảnh map đã xây dựng ở bài 04. Kiểm tra SHA256 xác nhận ảnh bài 04 và bài 05 giống nhau. YAML của bản sao quy định resolution $r=0.05$ m/pixel, origin $(-10.1,-10.2,0)$, occupied threshold 0.65 và free threshold 0.196. Ảnh map gốc không bị sửa để đặt thêm đồ vật.
\begin{lstlisting}[caption={Các trường của map YAML dùng trong bài 05},label={lst:map}]
image: map.pgm
mode: trinary
resolution: 0.05
origin: [-10.1, -10.2, 0.0]
negate: 0
occupied_thresh: 0.65
free_thresh: 0.196
\end{lstlisting}
Với ảnh xám $I\in[0,255]$ và negate bằng 0, xác suất chiếm chỗ được suy ra từ $p=1-I/255$. Pixel trắng là vùng trống, đen là vật cản; vùng nằm giữa hai ngưỡng là chưa biết. Ví dụ pixel 205 cho $p\approx0.196078>0.196$ nên giữ unknown, tránh coi vùng chưa quét là sàn trống.

Map server đọc metadata và ảnh, phát \code{/map} kiểu OccupancyGrid. Thông tin origin/resolution xác định phép chuyển tọa độ mét sang ô lưới. Với origin yaw bằng 0:
\begin{equation}
u=(x-x_0)/r,\qquad v=(y-y_0)/r.
\end{equation}
Khi truy cập ảnh PGM để vẽ, trục hàng ảnh đi xuống nên phải đảo hướng so với $y$ của map. Hàm \code{point_is_free()} làm việc trên OccupancyGrid và còn xử lý origin quay, biên map, unknown và ô có giá trị chiếm chỗ từ 50 trở lên.

\subsection{Bản đồ tĩnh và vật thể thêm vào simulator}
Các vật thể demo không được vẽ đè vào PGM để tạo kết quả localization giả. Static layer tiếp tục dùng bản đồ tuần trước. LiDAR quan sát vật thể mới và cập nhật obstacle/voxel layer của costmap. Inflation layer mở rộng chi phí xung quanh vật cản để planner/controller dành khoảng hở cho robot.

\note{Cùng một ảnh PGM chưa đủ bảo đảm cùng hệ tọa độ nếu đổi origin, resolution hoặc world. Các YAML waypoint/scene chứa hash ảnh; readiness còn kiểm tra điểm goal trên OccupancyGrid đang nhận. Scene chỉ được dùng với maze và map/world alignment đã xác thực.}
\clearpage

\subsection{AMCL: dự đoán, hiệu chỉnh và tái lấy mẫu}
AMCL biểu diễn phân bố pose robot bằng tập hạt $s_t^{(i)}=(x,y,\theta)$. Khi có chuyển động odometry $u_t$, hạt được dự đoán theo mô hình vi sai; khi có scan $z_t$, likelihood của từng hạt được tính bằng độ phù hợp giữa điểm LiDAR và tường trên map. Chu trình ở mức khái niệm:
\begin{align}
s_t^{(i)}&\sim p(s_t\mid s_{t-1}^{(i)},u_t),\\
\widetilde w_t^{(i)}&\propto p(z_t\mid s_t^{(i)},m),\qquad
w_t^{(i)}=\frac{\widetilde w_t^{(i)}}{\sum_j\widetilde w_t^{(j)}}.
\end{align}
Sau hiệu chỉnh, tái lấy mẫu tập trung hạt vào các pose có trọng số cao. Cấu hình đang dùng \code{likelihood_field}, \code{DifferentialMotionModel}, 500--2000 hạt, tối đa 60 beam và \code{resample_interval=1}. Các giá trị này đọc từ \code{turtlebot4_navigation/config/localization.yaml} của bản Humble cài trên máy.

\subsection{Chuỗi TF và pose khởi tạo}
\begin{figure}[H]\centering
\begin{tikzpicture}
\node[block,text width=3.4cm] (m) at (0,0) {\textbf{map}\\Hệ tọa độ toàn cục};
\node[block,text width=3.4cm] (o) at (5.1,0) {\textbf{odom}\\Liên tục cục bộ};
\node[block,text width=3.4cm] (b) at (10.2,0) {\textbf{base\_link}\\Thân robot};
\draw[arrow](m)--node[label,above]{AMCL}(o);
\draw[arrow](o)--node[label,above]{odometry}(b);
\end{tikzpicture}
\caption[Quan hệ các frame map, odom và base link]{Chuỗi TF dùng để xác định pose robot trong map. Các frame sensor được nối tiếp qua TF của robot.}\label{fig:tf}
\end{figure}
Odometry tích lũy chuyển động mượt nhưng trôi dần. AMCL ước lượng hiệu chỉnh map--odom dựa trên LiDAR; kết hợp hai biến đổi cho pose map--base\_link. Không nên có hai node cùng phát một biến đổi map--odom; vì vậy không chạy SLAM đồng thời với AMCL cho lượt demo này.

Trên RViz, chọn \textbf{2D Pose Estimate}, bấm vị trí robot và kéo mũi tên theo hướng đầu robot. Message \code{/initialpose} chứa pose và covariance để khởi tạo phân bố hạt. Sau đó kiểm tra scan khớp mép tường, localization active và TF hợp lệ. Việc thấy một mũi tên pose chưa tự chứng minh AMCL đúng.

Helper \code{initialize_pose} có chế độ \code{source:=sim_ground_truth}, dùng pose mô phỏng một lần trước mission khi hash map và alignment hợp lệ. Nếu dùng cách này phải nêu rõ trong video. Sau khởi tạo, định vị vẫn do AMCL thực hiện; ground truth chỉ phục vụ đối chiếu, không liên tục ghi đè AMCL khi robot chạy.
\clearpage

\section{Cách Nav2 lập kế hoạch và điều khiển robot}
\subsection{Từ mục tiêu tới vận tốc}
\begin{figure}[H]\centering
\begin{tikzpicture}[node distance=9mm]
\node[block,text width=11.5cm] (goal) {\textbf{NavigateToPose: PoseStamped trong frame map}};
\node[block,text width=11.5cm,below=of goal] (bt) {\textbf{BT navigator} -- điều phối lập đường, bám đường và recovery};
\node[block,text width=11.5cm,below=of bt] (planner) {\textbf{NavFn / Dijkstra + global costmap} $\rightarrow$ đường toàn cục};
\node[block,text width=11.5cm,below=of planner] (controller) {\textbf{DWB + local costmap + TF/odom} $\rightarrow$ vận tốc ứng viên tốt nhất};
\node[block,text width=11.5cm,below=of controller,fill=softgreen] (velocity) {\textbf{Velocity smoother $\rightarrow$ /cmd\_vel $\rightarrow$ Create 3/bridge}\\Robot chuyển động; scan, odom và TF tiếp tục cập nhật};
\draw[arrow](goal)--(bt);\draw[arrow](bt)--(planner);\draw[arrow](planner)--(controller);\draw[arrow](controller)--(velocity);
\end{tikzpicture}
\caption[Pipeline lập kế hoạch và điều khiển Nav2]{Chuỗi xử lý điều hướng của cấu hình đang dùng. Costmap và cảm biến được cập nhật trong suốt quá trình.}\label{fig:navflow}
\end{figure}
\textbf{Global planner.} Plugin \code{nav2_navfn_planner/NavfnPlanner}, id \code{GridBased}, có \code{use_astar=false}: sử dụng Dijkstra trên lưới costmap, lan truyền chi phí và truy hồi đường tới goal. Đây là cấu hình NavFn, không phải Smac hay thuật toán A* tự viết. Planner có tolerance 0.5 m và \code{allow_unknown=true}; dù vậy client vẫn yêu cầu vị trí goal nằm trong vùng known-free có margin.

\textbf{Local controller.} Plugin \code{dwb_core::DWBLocalPlanner} lấy mẫu vận tốc, mô phỏng quỹ đạo ngắn, loại quỹ đạo không hợp lệ và chấm điểm các ứng viên. Các critic gồm RotateToGoal, Oscillation, BaseObstacle, GoalAlign, PathAlign, PathDist, GoalDist. Ở mức khái niệm, lựa chọn tương ứng:
\begin{equation}
(v^*,\omega^*)=\arg\min_{(v,\omega)\in\mathcal V_{\rm valid}}\sum_k\alpha_k C_k(v,\omega).
\end{equation}
Công thức giải thích cơ chế chấm điểm, không thay thế mã DWB. Robot đi theo từng lệnh vận tốc ngắn và tính lại liên tục khi cảm biến/pose thay đổi.
\clearpage

\subsection{Costmap, thông số và Behavior Tree}
\begin{table}[H]\centering\small
\caption{Thông số Nav2 đọc từ cấu hình chính thức đang cài}\label{tab:navparams}
\begin{tabularx}{\linewidth}{@{}P{5.2cm}Y@{}}\toprule
Thành phần&Thông số và tác động\\\midrule
Global costmap&Frame map; resolution 0.06 m; static, obstacle, inflation layer; cập nhật 1 Hz; theo dõi unknown.\\
Local costmap&Frame odom; rolling window $3\times3$ m; resolution 0.06 m; static, voxel, inflation; cập nhật 5 Hz.\\
Inflation&Robot radius 0.175 m; inflation radius 0.45 m; cost scaling factor 4.0.\\
Controller&20 Hz; $v_x\le0.26$ m/s; $|\omega|\le1.0$ rad/s; DWB dự báo 1.7 s.\\
Goal checker&SimpleGoalChecker; xy tolerance 0.25 m; yaw tolerance 0.25 rad; stateful true.\\
Progress checker&Phải dịch chuyển 0.5 m trong thời gian cho phép 10 s theo cấu hình SimpleProgressChecker.\\
Velocity smoother&20 Hz; giới hạn vận tốc tuyến tính $[-0.26,0.26]$ m/s; góc $[-1,1]$ rad/s.\\\bottomrule
\end{tabularx}\end{table}
Các giá trị trên là cấu hình được launch kế thừa, không phải số đo tần số thực tế. Nếu mô phỏng có real-time factor thấp, thời gian thực để robot đi hết đường tăng lên dù thông số vận tốc trong thời gian mô phỏng giữ nguyên.

Behavior Tree mặc định \code{navigate_to_pose_w_replanning_and_recovery.xml} lập lại đường toàn cục ở nhánh RateController 1 Hz, đồng thời chạy FollowPath. Khi planner/controller lỗi, nhánh tương ứng có thể xóa costmap; recovery tổng thể lần lượt có clear costmap, spin, wait, backup. Không phải lúc nào cũng chạy tất cả recovery: chỉ nhánh được kích hoạt mới thực thi.

\subsection{Phân biệt các lớp điều phối}
BT quản lý cách đi đến \emph{một} goal. Node \code{sequential_navigation} quản lý \emph{danh sách} ba goal: gửi một action, chờ result, dwell rồi chuyển mục tiêu. Retry của node bài 05 và retry/recovery nội bộ BT là hai mức khác nhau; không cộng số lượng chúng để suy ra chính xác số lần robot quay/lùi.

Kết quả SUCCEEDED được Nav2 đánh giá trên pose ước lượng và goal checker. Nó không có nghĩa tọa độ ground truth trùng hoàn toàn với waypoint. Sai số đo sau action cũng chịu độ lệch thời điểm lấy mẫu của observer; báo cáo tách rõ hai nguồn pose khi đưa số liệu.
\clearpage

\section{Giải thích chương trình điều hướng tuần tự}
\subsection{Waypoint và kiểm tra trước khi gửi goal}
\begin{table}[H]\centering
\caption{Ba mục tiêu được đọc từ cấu hình waypoint}\label{tab:goals}
\begin{tabular}{@{}lrrr@{}}\toprule
Tên&$x$ (m)&$y$ (m)&yaw (rad)\\\midrule
Goal\_1&$-1.112140$&$0.315419$&$3.128185$\\
Goal\_2&$3.174917$&$-3.194671$&$3.128185$\\
Goal\_3&$3.615194$&$5.427635$&$0.258135$\\\bottomrule
\end{tabular}\end{table}
Tọa độ có nguồn từ lượt NavigateToPose thành công của bài 04, được lưu trong \code{waypoints.yaml} cùng source UUID. Goal\_3 sử dụng pose AMCL cuối đã đo, thay endpoint cũ quá sát vùng unknown. \code{load_waypoints()} bắt buộc validated, frame map, ít nhất ba goal, tên không trùng, source không rỗng, giá trị hữu hạn và hash ảnh map phù hợp.

Yaw không được truyền như một trường trực tiếp của Pose. Với chuyển động trên mặt phẳng, quaternion là:
\begin{equation}
(q_x,q_y,q_z,q_w)=\left(0,0,\sin\frac{\theta}{2},\cos\frac{\theta}{2}\right).
\end{equation}
\code{capture_waypoints} nhận \code{/clicked_point}, vốn chỉ có vị trí $x,y,z$, nên lấy yaw từ AMCL hoặc tham số được chọn riêng. Publish Point trên RViz không chứa hướng mũi robot.

\subsection{Readiness: vì sao chương trình có thể chờ?}
Trước goal đầu và giữa các goal, client kiểm tra: người dùng xác nhận localization; clock đang tiến và scan còn mới; đã nhận map và AMCL; covariance vị trí không quá 0.3, yaw không quá 0.2; TF map--base\_link không cũ quá 2 s; các lifecycle node active; action server có mặt; không có SLAM/teleop đang chạy; YOLO có inference thành công gần đây; toàn bộ vùng tròn quanh mỗi goal nằm trong known-free với margin 0.25 m.

Các giá trị covariance là phương sai (m$^2$, rad$^2$), không phải sai số mét/radian. Readiness không chứng minh localization đúng tuyệt đối: cần đối chiếu scan/map trên RViz trước khi đặt \code{localization_confirmed:=true}.
\note{YAML cấu hình navigation đặt dwell 8 s; default trong constructor là 6 s. Khi chạy đúng lệnh với \code{--params-file config/navigation.yaml}, giá trị có hiệu lực là 8 s. Không nên bỏ params-file rồi suy ra hành vi giống kịch bản quay.}
\clearpage

\subsection{Máy trạng thái, timeout và hủy goal}
\begin{figure}[H]\centering
\begin{tikzpicture}
\node[block,text width=3.5cm] (wait) at (0,0) {WAIT\_READY};
\node[block,text width=3.5cm] (send) at (0,-1.65) {SENDING};
\node[block,text width=3.5cm] (run) at (0,-3.3) {RUNNING};
\node[block,text width=3.5cm,fill=softgreen] (dwell) at (0,-4.95) {DWELL};
\node[block,text width=3.5cm,fill=softgreen] (done) at (0,-6.6) {DONE};
\node[block,text width=3.8cm,fill=softorange] (cancel) at (7,-3.3) {CANCELING};
\node[block,text width=3.8cm,fill=softorange] (retry) at (7,-1.65) {RETRY / FAILED};
\draw[arrow](wait)--node[label,right]{readiness đạt}(send);
\draw[arrow](send)--node[label,right]{server accepted}(run);
\draw[arrow](run)--node[label,right]{status = 4}(dwell);
\draw[arrow](dwell)--node[label,right]{đủ ba goal}(done);
\draw[arrow](run)--node[label,above]{timeout / Ctrl+C}(cancel);
\draw[arrow](cancel)--node[label,right]{kết thúc đã xác nhận}(retry);
\draw[arrow](retry.west)--node[label,above]{còn retry; chờ 2 s}(send.east);
\draw[arrow](dwell.west)--++(-1,0)|-node[label,pos=.72,left]{còn goal; chờ 8 s}(send.west);
\end{tikzpicture}
\caption[Máy trạng thái gửi goal tuần tự]{Luồng trạng thái rút gọn của client. Khi không xác nhận được hủy trong giới hạn 8 s, node dừng với lỗi và không gửi retry.}\label{fig:states}
\end{figure}
\begin{lstlisting}[language=Python,caption={Trích logic xử lý kết quả trong sequential\_navigation.py},label={lst:result}]
elif status == 4:
    self.index += 1
    self.attempt = 0
    self.state = 'DWELL'
    self.next_at = time.monotonic()+self.params['dwell_seconds']
    self.log('goal_succeeded', completed=self.index,
             dwell_seconds=self.params['dwell_seconds'])
else:
    self.failed(f'Action terminal status {status}; only 4 is success')
\end{lstlisting}
Readiness timeout 180 s; mỗi goal tối đa 300 s; tối đa một retry, chờ 2 s. Timeout dùng steady clock/wall monotonic để vẫn hết hạn khi simulator pause; header pose vẫn dùng ROS simulation time. Khi Ctrl+C, client gửi cancel và đợi kết quả, trả exit code 130. Nếu mất ACK/result, node không được tự suy ra rằng server đã dừng; nó báo lỗi sau giới hạn chờ, không gửi một goal khác có thể chạy chồng.

YOLO mất freshness khi đang chạy được đánh dấu \code{yolo_interrupted}; việc này không tự biến YOLO thành bộ lái hay tự làm action SUCCEEDED. Muốn chứng minh đủ yêu cầu video cần cả completed=3 và nhận diện chạy đồng thời, không chỉ nhìn dòng mission\_completed.
